%%%%%%%%%%%%%%%%%%%%%%%%%%%%%%%%%%%%%%%%%%%%%%%%%%%%%%%
% SECTION 6: RESULTS 2.3 - GPRC5D CASE STUDY
%%%%%%%%%%%%%%%%%%%%%%%%%%%%%%%%%%%%%%%%%%%%%%%%%%%%%%%
% 
% INPUTS NEEDED: 
%   - 15 resistance variants with details
%   - Three resistance categories
%   - Clonal fraction ranges
%
% FIGURES: None in main text (Supp Figure S1 in supplementary)
% TABLES: References Supplementary Table S2
% DEPENDENCIES: References [20, 21]
%
%%%%%%%%%%%%%%%%%%%%%%%%%%%%%%%%%%%%%%%%%%%%%%%%%%%%%%%

\subsection{Case study: GPRC5D in relapsed/refractory myeloma}\label{subsec3}

%%% PARAGRAPH 1: Rationale for GPRC5D as validation target %%%
To validate the clinical utility of our ReACT-enhanced extraction pipeline for capturing complex resistance mechanisms in emerging therapeutics, we applied the system to process literature describing GPRC5D-targeted therapy resistance in multiple myeloma (MM), a cancer of plasma cells. GPRC5D represents an ideal test case as an emerging target for CAR T-cell and bispecific antibody therapies with limited existing database coverage despite active clinical deployment. Current CIViC entries for GPRC5D are minimal, yet recent literature documents resistance patterns that directly impact therapeutic decision-making and patient outcomes.

%%% PARAGRAPH 2: Pipeline application and capabilities %%%
Our pipeline processed detailed case studies of GPRC5D-targeted therapy resistance, systematically identifying and structuring discrete molecular escape mechanisms across multiple patient cases\cite{bib20,bib21}. The ReACT reasoning framework supported structured clinical interpretation by recognizing complex polyclonal resistance evolution, distinguishing between genetic, epigenetic, and structural mechanisms, and maintaining precise quantitative tracking of clonal dynamics throughout therapeutic courses.

%%% PARAGRAPH 3: Evidence item generation summary %%%
The extraction pipeline successfully generated 15 CIViC-compatible evidence items from the recent literature\cite{bib20,bib21}, each representing a distinct molecular mechanism of therapeutic escape. Every evidence item adheres to CIViC curation standards with appropriate Evidence Type (Predictive), Evidence Level (C for case studies), Evidence Direction (Supports), and Clinical Significance (Resistance). The comprehensive Evidence Statements include patient identifiers, quantitative clonal fractions, clinical timelines, and mechanistic details essential for clinical application.

%%% PARAGRAPH 4: Three resistance categories breakdown %%%
The extracted evidence encompasses three major resistance categories: genetic mechanisms (9 evidence items including point mutations Arg233Ter, Tyr257Ser, Glu146Ter), frameshift mutations (E27fs, S125fs, F158fs, Leu174TrpfsTer180), and in-frame deletions (G97\_F100del) with clonal fractions ranging from 6.2\% to 45\%; epigenetic mechanisms (1 evidence item) involving promoter hypermethylation resulting in greater than 99.8\% cells lacking GPRC5D expression; and structural mechanisms (4 evidence items) including approximately 220~kb focal deletions, approximately 267~kb biallelic deletions, and chromosome 12p arm deletions demonstrating convergent evolution toward complete target antigen loss.

%%% PARAGRAPH 5: Summary statement %%%
This case study demonstrates the system's ability to identify and structure emerging therapeutic evidence that would otherwise remain buried in narrative reviews, enabling rapid database updates for novel targets entering clinical practice. The complete GPRC5D evidence table with all 15 entries is provided in Supplementary Table~S2.

%%%%%%%%%%%%%%%%%%%%%%%%%%%%%%%%%%%%%%%%%%%%%%%%%%%%%%%
% KEY DATA IN THIS SECTION:
%
% TOTAL: 15 CIViC-compatible evidence items
%
% THREE RESISTANCE CATEGORIES:
%
% 1. GENETIC MECHANISMS (9 items):
%    Point mutations:
%      □ Arg233Ter
%      □ Tyr257Ser
%      □ Glu146Ter
%    Frameshift mutations:
%      □ E27fs
%      □ S125fs
%      □ F158fs
%      □ Leu174TrpfsTer180
%    In-frame deletions:
%      □ G97_F100del
%    Clonal fractions: 6.2% to 45%
%
% 2. EPIGENETIC MECHANISMS (1 item):
%      □ Promoter hypermethylation
%      □ >99.8% cells lacking GPRC5D expression
%
% 3. STRUCTURAL MECHANISMS (4 items):
%      □ ~220 kb focal deletion
%      □ ~267 kb biallelic deletion
%      □ Chromosome 12p arm deletion
%      □ Convergent evolution toward target antigen loss
%
% CIVIC STANDARDS APPLIED:
%      □ Evidence Type: Predictive
%      □ Evidence Level: C (case studies)
%      □ Evidence Direction: Supports
%      □ Clinical Significance: Resistance
%
%%%%%%%%%%%%%%%%%%%%%%%%%%%%%%%%%%%%%%%%%%%%%%%%%%%%%%%

%%%%%%%%%%%%%%%%%%%%%%%%%%%%%%%%%%%%%%%%%%%%%%%%%%%%%%%
% SUPPLEMENTARY MATERIALS REFERENCED:
%
% → Supplementary Table S2: Complete GPRC5D evidence table (15 entries)
%   Contains: Variant, Type, Patient, Therapy, Clonal %, Timeline
%
% → Supplementary Table S3: Detailed clinical evidence statements
%   Contains: Full evidence statements with DOIs for all 15 variants
%
% → Supplementary Figure S1: GPRC5D case study workflow
%   Shows: ReACT extraction workflow for resistance mechanisms
%
%%%%%%%%%%%%%%%%%%%%%%%%%%%%%%%%%%%%%%%%%%%%%%%%%%%%%%%

%%%%%%%%%%%%%%%%%%%%%%%%%%%%%%%%%%%%%%%%%%%%%%%%%%%%%%%
% REFERENCES USED:
%   [20] Lee et al. 2023 - Nature Medicine
%        "Mechanisms of antigen escape from BCMA- or GPRC5D-
%        targeted immunotherapies in multiple myeloma"
%        DOI: 10.1038/s41591-023-02491-5
%
%   [21] Derrien et al. 2023 - Nature Cancer
%        "Acquired resistance to a GPRC5D-directed T-cell 
%        engager in multiple myeloma is mediated by genetic 
%        or epigenetic target inactivation"
%        DOI: 10.1038/s43018-023-00625-9
%%%%%%%%%%%%%%%%%%%%%%%%%%%%%%%%%%%%%%%%%%%%%%%%%%%%%%%
